\chapter{Results}

\section{Neural Network Approximation}

The final neural network results on the held-out test set are reported in Table~\ref{tab:nn-final-metrics}.

\begin{table}[H]
    \centering
    \begin{tabular}{lr}
        \toprule
        Metric & Value \\
        \midrule
        MAE & 0.04290 \\
        RMSE & 0.06208 \\
        $R^2$ & 0.999993 \\
        MAPE, prices $> 1$ & 0.4803\% \\
        \bottomrule
    \end{tabular}
    \caption{Final neural network metrics against analytical Black-Scholes prices.}
    \label{tab:nn-final-metrics}
\end{table}

The results show that the feed-forward neural network with moneyness as an engineered feature and SiLU hidden activations approximates the Black-Scholes pricing function with very high accuracy on the sampled parameter domain.
The coefficient of determination is close to one, while MAE and RMSE are small in absolute price units.
This supports the main research question: in this controlled setting, a neural network can learn the Black-Scholes pricing map effectively and can be interpreted as a fast pricing surrogate after training.

\begin{figure}[H]
    \centering
    \maybeincludegraphics[width=0.72\linewidth]{../results/final/figures/loss.png}
    \caption{Training loss and validation loss for the final experiment.}
    \label{fig:final-loss}
\end{figure}

Figure~\ref{fig:final-loss} shows that the optimization process remains stable.
The validation curve follows the training curve closely, suggesting that the model does not simply memorize the training set.

\begin{figure}[H]
    \centering
    \maybeincludegraphics[width=0.62\linewidth]{../results/final/figures/true_vs_predicted.png}
    \caption{Analytical Black-Scholes prices against neural network predictions.}
    \label{fig:final-true-vs-predicted}
\end{figure}

Figure~\ref{fig:final-true-vs-predicted} confirms that predicted prices are tightly aligned with analytical prices across the observed price range.

\section{Error Diagnostics}

\begin{figure}[H]
    \centering
    \maybeincludegraphics[width=0.72\linewidth]{../results/final/figures/error_distribution.png}
    \caption{Distribution of neural network prediction errors.}
    \label{fig:final-error-distribution}
\end{figure}

The error distribution in Figure~\ref{fig:final-error-distribution} is concentrated near zero, which is consistent with the low MAE and RMSE values reported in Table~\ref{tab:nn-final-metrics}.

\begin{figure}[H]
    \centering
    \begin{subfigure}{0.32\linewidth}
        \centering
        \maybeincludegraphics[width=\linewidth]{../results/final/figures/error_vs_moneyness.png}
        \caption{Moneyness}
    \end{subfigure}
    \begin{subfigure}{0.32\linewidth}
        \centering
        \maybeincludegraphics[width=\linewidth]{../results/final/figures/error_vs_maturity.png}
        \caption{Maturity}
    \end{subfigure}
    \begin{subfigure}{0.32\linewidth}
        \centering
        \maybeincludegraphics[width=\linewidth]{../results/final/figures/error_vs_volatility.png}
        \caption{Volatility}
    \end{subfigure}
    \caption{Absolute neural network error against selected financial features.}
    \label{fig:final-error-diagnostics}
\end{figure}

Figure~\ref{fig:final-error-diagnostics} is used to inspect whether the approximation error changes systematically across the parameter space.
This is important because aggregate metrics alone may hide localized weaknesses, especially around different moneyness regimes, short maturities, or high-volatility regions.

\section{Monte Carlo Benchmark}

The Monte Carlo benchmark is evaluated against the analytical Black-Scholes price on 512 deterministic test options.
The results are shown in Table~\ref{tab:mc-final-metrics}.

\begin{table}[H]
    \centering
    \begin{tabular}{lr}
        \toprule
        Metric & Value \\
        \midrule
        MAE & 0.08882 \\
        RMSE & 0.15790 \\
        $R^2$ & 0.999951 \\
        MAPE, prices $> 1$ & 0.6483\% \\
        \bottomrule
    \end{tabular}
    \caption{Final Monte Carlo metrics against analytical Black-Scholes prices.}
    \label{tab:mc-final-metrics}
\end{table}

The Monte Carlo benchmark also achieves high accuracy, as expected when the number of simulated paths is large.
However, it remains a simulation-based estimator and therefore has sampling error.
The neural network, in contrast, produces deterministic forward-pass predictions once trained.
This highlights the main computational role of the network: it does not replace the analytical Black-Scholes formula in this setting, but it demonstrates how a learned surrogate can approximate a pricing map efficiently.

\section{Runtime Benchmark}

The runtime benchmark separates the initial neural network training cost from pricing time after training.
Using the improved final configuration, training required approximately 359 seconds in the benchmark environment.
Once trained, the neural network priced 1,000 options in approximately 0.0319 seconds on average.
The Monte Carlo benchmark with 20,000 paths required approximately 1.2824 seconds for the same number of options, while the vectorized analytical Black-Scholes formula required approximately 0.00091 seconds.

\begin{table}[H]
    \centering
    \begin{tabular}{lrr}
        \toprule
        Method & Mean time for 1,000 options & Mean options/s \\
        \midrule
        Black-Scholes formula & 0.00091 s & 1,101,094.82 \\
        Neural network inference & 0.03190 s & 31,350.77 \\
        Monte Carlo, 20,000 paths & 1.28235 s & 779.82 \\
        \bottomrule
    \end{tabular}
    \caption{Runtime benchmark for analytical, neural-network, and Monte Carlo pricing.}
    \label{tab:runtime-benchmark}
\end{table}

As expected, the analytical Black-Scholes formula is the fastest method when available.
The relevant surrogate-pricing comparison is therefore between neural network inference and Monte Carlo simulation.
In this benchmark, neural inference is roughly 40 times faster than Monte Carlo after the model has been trained.

\section{Discussion}

The final results indicate that the neural network performs better than the Monte Carlo benchmark in terms of MAE, RMSE, and MAPE on the reported experiment.
This comparison should be interpreted carefully because the two methods are evaluated differently: the neural network is evaluated on the full held-out test set, whereas Monte Carlo is evaluated on a deterministic subset due to computational cost.

The most important conclusion is qualitative rather than competitive.
In a controlled Black-Scholes environment, a feed-forward network can learn the mapping from option and market parameters to theoretical prices with high accuracy.
This validates the use of the Black-Scholes setting as a clean benchmark for studying neural-network pricing surrogates before moving to more complex models.
